\chapter{Results and Comparison}
\label{ch:ch7_results}

This chapter collects what the work establishes and separates it, without
softening, from what it does not. The engine has been modelled, written,
verified in simulation against official vectors and synthesized on a real
standard-cell library; it has not yet run on the \ac{FPGA}. Every number
below therefore carries its provenance tag, and the entropy campaign ---
the one result that requires silicon on the bench --- is stated as a
protocol to execute, not as data.

\section{Status against the objectives}
\label{sec:status}

Table~\ref{tab:status} maps the objectives of
Section~\ref{sec:objectives} onto evidence.

\begin{table}[htbp]
\centering
\caption{Objectives and the evidence supporting each.}
\label{tab:status}
\small
\begin{tabular}{@{}llp{7.4cm}@{}}
\toprule
Obj. & State & Evidence \\
\midrule
O1 & met    & Ring described at transistor level; noise swept in SPICE
              (Section~\ref{sec:spice_results}) \\
O2 & met    & Jitter extracted from the waveforms; estimator validated on
              synthetic data with known $Q$ \\
O3 & met    & Stochastic model closed: exact wrapped-Gaussian computation,
              design point $Q \ge 0.2286$ (Section~\ref{sec:opoint}) \\
O4 & met    & Engine written in portable VHDL-2008; seven testbenches pass
              against FIPS-197, SP 800-38B, RFC 4493 and \ac{CAVP} vectors \\
O5 & partial& Verifier firmware compiles and the protocol is closed; the
              board-level run is pending \\
O6 & pending& Requires the \ac{FPGA} on the bench
              (Section~\ref{sec:protocol}) \\
O7 & met    & Synthesis on IHP SG13G2, per-block area and optimization path
              (Chapter~\ref{ch:ch6_asic}) \\
O8 & partial& Structural comparison closed below; the numerical comparison
              against the ESP32 needs O6 \\
\bottomrule
\end{tabular}
\end{table}

The minimum defensible result declared in Chapter~\ref{ch:ch1_intro} ---
O1--O3 together with O5 --- is reached on its modelling and design side.
What O5 still lacks is the bring-up, not the design.

\section{Transistor-level noise of the ring}
\label{sec:spice_results}

The ring was simulated with a current noise source of controlled spectral
density injected at the switching node, sweeping the density and measuring
the standard deviation of the rising-edge crossings over 500 periods.
Table~\ref{tab:spice} gives the result for a five-stage ring with a
10~fF load.

\begin{table}[htbp]
\centering
\caption{Period jitter of the five-stage ring against injected noise
density \simu{}. $\sigma_{\mathrm{rel}} = \sigma/T$.}
\label{tab:spice}
\small
\begin{tabular}{@{}lrrrr@{}}
\toprule
Noise density (A$^2$/Hz) & Crossings & $T$ (ps) & $\sigma$ (fs) &
$\sigma_{\mathrm{rel}}$ \\
\midrule
none (numerical floor) & 500 & 230.30 & 1.38   & $6.0\times10^{-6}$ \\
$1\times10^{-20}$      & 501 & 230.37 & 1200   & $5.2\times10^{-3}$ \\
$4\times10^{-20}$      & 507 & 229.98 & 2416   & $1.1\times10^{-2}$ \\
$1.6\times10^{-19}$    & 509 & 227.77 & 25588  & $1.1\times10^{-1}$ \\
\bottomrule
\end{tabular}
\end{table}

Three readings, and the third is the one that matters.

\paragraph{The numerical floor is far enough below the signal.} With the
noise source switched off the measurement returns 1.38~fs, which is the
resolution of the simulation itself, not physics. It sits about 870 times
below the smallest jitter to be measured, so the floor does not
contaminate the sweep. Reaching it required care: the first attempt
returned 0.33~ps, of the same order as the quantity under study, and the
cause was the ring's start-up transient being included in the statistics.
Discarding the first twenty cycles removed it. A measurement floor that
coincides with the signal is indistinguishable from a result, which is
precisely why the floor was characterized before the sweep and not after.

\paragraph{The scaling law holds where it should.} Thermal jitter
accumulates as the square root of the noise density, so quadrupling the
density must double $\sigma$. From $1\times10^{-20}$ to
$4\times10^{-20}$ the measurement gives a factor of 2.01 against a
predicted 2.00. That agreement is the evidence that what is being measured
is the injected noise and not an artefact of the solver.

\paragraph{And it breaks at the top of the sweep, visibly.} At
$1.6\times10^{-19}$, sixteen times the reference density, $\sigma$ grows by
a factor of 21.3 where the law predicts 4.0, and the mean period drops from
230.4 to 227.8~ps. A zero-mean noise source cannot shift a mean period.
The point is therefore outside the regime the model describes --- the
perturbation is large enough to change the operating point of the
inverters --- and it is reported rather than dropped, because the boundary
of validity of a model is itself a result. Only the two central points are
used for anything downstream.

\section{Operating point: divider and throughput}
\label{sec:opoint}

The stochastic model fixes the target. For the PTG.2 thresholds of AIS
20/31 v3.0, the min-entropy criterion ($H_\infty \ge 0.98$) demands
$Q \ge 0.2286$, and is more demanding than the Shannon criterion
($H \ge 0.9998$, $Q \ge 0.2022$); the design therefore targets the former,
with $Q = 0.30$ as the working set-point, which yields
$H_\infty = 0.9951$. Two properties of the model were established rather
than assumed: the truncated series of equation~14
in~\cite{baudet_2011_security} \emph{overestimates} entropy and is thus
not a lower bound, so the exact wrapped-Gaussian computation is used
instead; and the first-harmonic derivation of min-entropy agrees with the
exact one to six digits over the whole range of interest.

Since $Q_{\mathrm{gen}} = K_D \cdot Q_{\mathrm{raw}}$, the divider follows
from the jitter that the campaign measures. Table~\ref{tab:opoint} gives
the operating points for the jitter range expected of this
technology~\cite{petura_2016_survey}, assuming RO2 at 500~MHz.

\begin{table}[htbp]
\centering
\caption{Divider and resulting throughput per ring for
$Q \ge 0.2286$ \est{}.}
\label{tab:opoint}
\small
\begin{tabular}{@{}rrrr@{}}
\toprule
$\sigma/T$ per period & $Q_{\mathrm{raw}}$ & $K_D$ & Throughput \\
\midrule
$7.1\times10^{-4}$ & $5\times10^{-7}$ & 457\,200 & 1.09~kbit/s \\
$1.4\times10^{-3}$ & $2\times10^{-6}$ & 114\,300 & 4.37~kbit/s \\
$2.8\times10^{-3}$ & $8\times10^{-6}$ & 28\,575  & 17.5~kbit/s \\
\bottomrule
\end{tabular}
\end{table}

These bracket the published design point of Fischer and
Lubicz~\cite{fischer_2014_embedded} ($K_D \approx 430\,000$ for
$\sigma = 5.01$~ps on $T = 7.81$~ns), which is the sanity check that the
model is being applied as its authors intended.

The architectural consequence is stated plainly: a single \ac{ERO}
delivers kbit/s, not Mbit/s. That is the price of demonstrable entropy,
and it is affordable here because the engine's product is a 256-bit seed
for the \ac{DRBG}, not a bit stream. More throughput means more
independent rings, each with its own bound, not a faster ring.

\section{Functional verification, consolidated}
\label{sec:verif_results}

Chapter~\ref{ch:ch5_verification} describes the strategy; this section
records the totals. Every figure in Table~\ref{tab:verif} was obtained by
running the suite, not by reading a previous log.

\begin{table}[htbp]
\centering
\caption{Verification totals \simu{}.}
\label{tab:verif}
\small
\begin{tabular}{@{}lll@{}}
\toprule
Suite & What it checks & Result \\
\midrule
VHDL testbenches (GHDL) & rings, \ac{ERO}, capture, health tests,
  \ac{I2C}, \ac{AES}, \ac{CMAC}, \ac{CTRDRBG}, top level & 7/7 pass \\
\texttt{ctr\_drbg\_ref.py} & \ac{CAVP} vectors, with and without
  derivation function & 960/960 \\
\texttt{ero\_model.py} & stochastic model, exactness and bounds & 24/24 \\
\texttt{jitter\_estimator.py} & jitter recovery on synthetic data of
  known $Q$ & 20/20 \\
\texttt{cortes\_salud.py} & health-test cutoffs against SP 800-90B
  formulas & 6/6 \\
\texttt{verifica\_top.py} & cross-check of the top-level bench with an
  independent \ac{AES} & 5/5 \\
\bottomrule
\end{tabular}
\end{table}

Two of these deserve a sentence. The 960 \ac{CAVP} cases are the
strongest single piece of evidence in the work: the generator is not
merely self-consistent, it reproduces the outputs that \ac{NIST} publishes
for \ac{CTRDRBG} with AES-128, in both the derivation-function and the
raw-entropy configurations. And the health-test cutoffs, $C_{RCT} = 22$
and $C_{APT} = 596$, had been computed by hand from the standard and were
never checked; deriving them again from the formulas of SP
800-90B~\cite{turan_2018_sp80090b} confirms both. The same check
quantifies the cost of the complement test added to the \ac{APT}, which is
an addition of this work and not of the standard: it raises the
false-alarm rate per window from $2^{-20.15}$ to $2^{-20.14}$, a 0.8\,\%
increase, because the complementary tail lies 5.7 standard deviations
away. The cutoff therefore remains valid as it stands and the addition is
free in practice.

\section{Comparison with commercial secure elements}
\label{sec:se_comparison}

The engine occupies the same functional slot as a commercial secure
element: a small companion device that generates and holds keys for a host
microcontroller over a two-wire bus. Table~\ref{tab:se} places it beside
the three parts that dominate the \ac{I2C} secure-element market.

\begin{table}[htbp]
\centering
\caption{The engine against commercial secure elements. Certification
data from the manufacturers' datasheets.}
\label{tab:se}
\small
\begin{tabular}{@{}lllll@{}}
\toprule
 & ATECC608B & SE050 & OPTIGA Trust M & This work \\
 & \cite{microchip_atecc608b_ds} & \cite{nxp_se050_ds} &
   \cite{infineon_optiga_trustm_ds} & \\
\midrule
Bus            & \ac{I2C} & \ac{I2C} & \ac{I2C} & \ac{I2C} \\
Certification  & FIPS 140-3 L2 & CC EAL6+ & CC EAL6+ & none \\
Entropy source & undisclosed & undisclosed & undisclosed & \ac{ERO},
                 published \\
Entropy bound  & not published & not published & not published &
                 derived from measured jitter \\
Key export     & not possible & not possible & not possible &
                 test mode only \\
Inspectable    & no & no & no & yes, \ac{RTL} released \\
\bottomrule
\end{tabular}
\end{table}

The honest reading of that table has two halves.

On capability and assurance the commercial parts win, and it is not close.
Two of them carry CC EAL6+ evaluations, which represent an amount of
third-party scrutiny --- including physical attack resistance --- that no
academic project reproduces. They offer asymmetric cryptography, certified
provisioning flows and manufacturing support. This work offers symmetric
key generation and authentication, and carries no certification at all.

On what the thesis is actually about, the comparison inverts. In all three
commercial devices the entropy source is undisclosed: the datasheets state
that a \ac{TRNG} is present, not what it is, and no min-entropy bound
derived from a measured physical parameter is published for any of them.
Their randomness is trusted because the vendor and the evaluation scheme
are trusted. The engine described here makes the opposite trade: it has no
certificate, and it can show where each bit of entropy comes from, what
physical measurement bounds it, and what the bound becomes when the
temperature or the divider changes. That is the same distinction that the
author's earlier characterization of the ESP32 generator ran into from the
other side --- a black box can be tested statistically but cannot be
bounded --- and it is the reason this engine was built as a white box.

Which of the two is preferable depends on the question being asked. For
shipping a product, the certified part. For knowing what the randomness in
that product is worth, only the open one answers.

\section{Uncertainty budget}
\label{sec:uncertainty}

Four quantities carry the conclusions, and each has a different and
declared kind of uncertainty.

\begin{description}
\item[Simulated jitter.] Resolution floor 1.38~fs \simu{}, about 870 times
below the smallest quantity measured. The scaling check against the
expected square-root law agrees to 0.5\,\% in the valid regime, and the
regime boundary is documented in Section~\ref{sec:spice_results}. The
figure depends on the noise density assumed for the technology, which is
itself \est{}: what the simulation establishes is the \emph{method} and
the ring's response, not the absolute jitter of a fabricated device.
\item[Recovered $Q$.] The estimator reproduces a known $Q$ with a bias
below 2\,\% and a dispersion of 4.0\,\% over eight seeds with 20\,000
blocks \simu{}. Without the corrections applied here a plain straight-line
fit biases the result by $-18$\,\%, so the correction is not cosmetic.
\item[Area.] A lower bound \simu{}: cells only, flattened, without routing,
fill or clock tree. Published growth factors from cell area to core area
on this library are $\approx 2.35$ for an \ac{AES}, and the optimization
savings of Chapter~\ref{ch:ch6_asic} are \est{} and explicitly not claimed
as measured.
\item[Min-entropy of the built device.] Not yet bounded. The model, the
estimator and the cutoffs are ready; the number requires the campaign of
Section~\ref{sec:protocol}. No value is offered here, and none should be
inferred from the simulated jitter.
\end{description}

\section{The pending campaign}
\label{sec:protocol}

The one result that silicon must supply is the min-entropy of the rings as
built, and the procedure is fixed in advance so that it cannot be tuned
after seeing the data.

\begin{enumerate}
\item \textbf{Confirm the rings oscillate.} Read the frequency of each
ring with the embedded counter before trusting any entropy measurement. A
ring removed by the synthesizer, or two rings that have locked to each
other, would otherwise be indistinguishable from a working source.
\item \textbf{Measure the jitter} with $K_D = 1$ by the method of Fischer
and Lubicz~\cite{fischer_2014_embedded}, sweeping the pair distance $M$
and fitting the linear region, with at least three rings so that the
global component separates from the local one~\cite{lubicz_2024_jitter}.
Fit $\sigma^2(t) = at + bt^2$ and keep only the thermal term.
\item \textbf{Fix $K_D$} from the measured jitter for $Q \ge 0.2286$ with
margin, and recompute the health-test cutoffs for the measured $H_\infty$
rather than the assumed 0.98, using \texttt{cortes\_salud.py}.
\item \textbf{Collect raw bits} at that operating point and run the
\ac{NIST} SP 800-90B estimators, contrasting the empirical estimate with
the model's bound. Agreement supports the model; disagreement is a result
about the model's hypotheses, and either way it is reported.
\item \textbf{Repeat across placement} --- rings in separated regions
versus adjacent ones --- to quantify the coupling that the literature
warns about, which is the measurement the architecture of this engine was
chosen to make possible.
\end{enumerate}

Until step 4 closes, the engine's claim is precise and limited: the chain
from physical noise to key material is designed, implemented and verified
against the standards, and the bound it will produce is computable from a
measurement that has not yet been taken.
